\summarychapter{12}{Exiles from Life}{The Beginnings of Monasticism}

\begin{inbrief}
When the empire turned Christian, the hero of the faith was no longer the martyr in the arena but the
hermit in the desert. The chapter follows the monastic movement from Anthony and the Egyptian hermits
through the common life of Pachomius and Basil to the Rule of Benedict, and weighs what the monks gave
the church against the flaw in their view of human life.
\question{How does the monks' renunciation relate to the gospel: is it salvation by works, or a
legitimate form of repentance?}
\end{inbrief}

Athanasius's \emph{Life of Saint Anthony} tells how Satan sent the wild beasts of the desert against
the praying monk, who told them to attack if God permitted and otherwise to be gone, ``for I am a
servant of Jesus the Conqueror''; Satan fled like a sparrow before a hawk. That was the new ideal.
Moderns, like Gibbon scoffing at ``unhappy exiles from social life,'' hardly know what to make of it,
and Christians disagree. Roman Catholics hold that the church has room both for ascetics who strive
for perfection and for weak members. The Reformers, led by the former monk Luther, objected that
monasticism teaches two roads to God, a higher and a lower, whereas the gospel knows one way, faith
in Christ active in love. All agree that the monks were ascetics who believed that renouncing the
body frees the soul for God.

\ssection{The Monastic Ideal}
Ascetic notes sound in John the Baptist, in Jesus' word to the rich young man and in Paul's conflict
of flesh and Spirit. About 140 the \emph{Shepherd of Hermas} already distinguished the precepts
binding on all from advice for those who want to do more. Praise of self-denial, especially celibacy,
grew; the practice of penance encouraged acts of exceptional virtue; Tertullian, Origen and Cyprian
supported a higher level of sanctity, and long before Constantine many took vows of abstinence while
still living in the cities.

The first monks were hermits (from the Greek for desert), and the flight began in Egypt, where the
desert lies close on either side of the Nile. Anthony, born about 250 in the village of Koma, gave
away his wealth at twenty in obedience to Christ's word to the rich young ruler and lived in solitude,
at first in a tomb, battling temptations that legend pictures as devils, beasts and women; he is said
to have lived 105 years. He had hundreds of imitators. The surge coincided with the sudden popularity
of Christianity. Constantine's conversion lowered commitment: half-converted pagans replaced the
martyrs, and Christians now fought each other for the prizes of the church. The hermit fled not so
much the world as the world in the church. Against the great institution he set the soul face to face
with God, which exposed him to individualism and to the inner temptations of pride, rivalry and
eccentricity. In Egypt and Syria some ate only grass, lived in trees or never washed. Simeon
Stylites, to escape the crowds, lived more than thirty years on top of a pillar, from which he
preached and converted thousands.

\ssection{The Move to Community Life}
About 320 Pachomius, a former soldier, founded the first monastery. In place of hermits who were each
a law to himself he set up a regulated common life (\emph{coenobitic} monasticism, from \emph{koinos
bios}) with fixed hours, manual work, uniform dress and strict discipline. It checked idleness and
eccentricity and opened the monastic life to women. ``To save souls you must bring them together.''
From Egypt the movement spread to Syria, Asia Minor and western Europe. Basil (died 379), with the two
Gregories, won Asia Minor and wrote the Rule under which Greek Orthodox monasticism is still
organized.

In the fourth and fifth centuries the ideal seized all levels of the Christian population. It
replaced the heroism of the age of persecution, revived the enthusiasm and end-time piety of earlier
days, and turned the spirit of martyrdom into total commitment and the ascetic imitation of Christ.
To live for God alone the monks took a threefold vow of poverty, chastity and obedience, giving up
what we call basic rights --- property, marriage, freedom of choice --- as roots of the ego. Once the
early extravagances passed and livable rules prevailed, monasteries took on tasks of great benefit.
In the fifth and sixth centuries almost every church leader was a monk or close to the monks, and the
cell became a study. The pioneer of monastic scholarship was Jerome (340--420), who began as a hermit
in the Syrian desert, mastered his temptations by learning Hebrew, taught Bishop Damasus and a circle
of noble ladies in Rome, and finally withdrew to Bethlehem, where he translated the Bible from the
original languages into Latin: the Vulgate, the authorized Bible of Roman Catholicism until recent
years. Athanasius brought monasticism to the West when he was banished to Trier in 335 with two monks,
and his \emph{Life of Anthony} spread the idea. Ambrose and Augustine were eager supporters; Augustine
wrote the first western rule for his clergy at Tagaste and Hippo, and in 415 John Cassian, who had
visited Egypt, founded St.~Victor near Marseille.

\ssection{The Genius of the West}
The constitution of Western monasticism came from Benedict, born at Nursia, northeast of Rome, late in
the fifth century. He broke off his studies in Rome to live three years as a hermit in a cave. Monks
nearby chose him as their abbot, found his discipline irksome and tried to poison him. In 529 he
founded on Monte Cassino the most famous monastery in Europe, the motherhouse of the Benedictine
order, wrote his Rule for it and lived there until his death in 542. He was no scholar, says Williston
Walker, but he had the Roman gift for administration, a firm belief in monasticism and a deep
knowledge of people. Seeing that monasticism's evils came from lack of discipline, he made discipline
fundamental yet bearable for ordinary men; this combination of restraint with some freedom is the
mark of his Rule.

He thought of the monastery as a garrison of Christ in a hostile world. A candidate tried the life for
at least a year before taking the lifelong vows of poverty, chastity and obedience. The abbot was owed
absolute obedience, but he was elected by the monks and had to consult the whole body on major
matters. Each house was to supply all its own needs --- cloth, wine, carpentry, masonry --- since
wandering outside was a danger to the soul. Worship filled the day: seven services, brief and made up
mostly of Psalms, and a vigil at two in the morning. Most fruitful was the rule on work: ``Idleness
is hostile to the soul,'' so the brothers spent fixed hours in manual labor and in religious reading,
more in the fields in summer, more with books in winter and Lent. A true Benedictine house was a
little world of worship, labor and study with its own library, and the monks soon copied the Latin
classics as well as the church fathers, which is why these survived.

\ssection{Some Pros and Cons}
The Rule spread quickly over western Europe, and the service of the monks after the ruin of Roman
civilization can hardly be exaggerated. Their houses gave the Middle Ages their only chance for
study, protection in constant war and rest; they were a great missionary force and a reminder to a
rough population that man does not live by bread alone. The faults are equally plain. While monks
vowed poverty, monasteries grew rich in land; discipline slackened, and medieval history is full of
reforms and new foundations meant to correct the old. Above all, Shelley judges with Walker, the
monastic idea of the Christian life was unnatural: it abandoned ordinary social relationships,
marriage and the home, and rested on a false view of man, the soul chained to the flesh like a
prisoner to a corpse, which is not the Bible's view. Those faults were not visible to the people of
the declining empire and the Middle Ages, for whom the monk's calling seemed the truest Christian
life, and they should not lead us to underrate what the monks did for Christianity and civilization
in a hard age.

\begin{keyfacts}
\keypeople{Anthony of Egypt; Athanasius; Simeon Stylites; Pachomius; Basil; Jerome; Damasus; Ambrose;
Augustine; John Cassian; Benedict of Nursia.}
\keydates{%
  \kd{c. 140}{The \emph{Shepherd of Hermas} speaks of a higher and a lower morality}
  \kd{c. 250}{Birth of Anthony, regarded as the first monk}
  \kd{c. 320}{Pachomius founds the first monastery with a common rule}
  \kd{335}{Athanasius, exiled to Trier, brings monasticism to the West}
  \kd{379}{Death of Basil, author of the Rule of Eastern monasticism}
  \kd{415}{John Cassian founds the monastery of St.~Victor near Marseille}
  \kd{529}{Benedict founds Monte Cassino and writes his Rule}}
\keyterms{Asceticism; hermit; celibacy; coenobitic (common) life; poverty, chastity, obedience; abbot;
Rule of Benedict; vigil; Vulgate.}
\end{keyfacts}
